\begin{abstract}
Language models can apply qualitative codebooks, but their labels save effort only if a researcher knows which ones to trust.
We ask whether model confidence can tell, comparing seven decision models (APIs that return probabilities for typed
questions) with twelve LLM confidence signals on five human-coded tasks spanning Reddit comments, counselling dialogue and
classroom talk, under a prospectively specified analysis plan. Three findings stand out. First, low calibration error does
not make confidence useful: decision models differ widely in how well they rank their own mistakes, and calibration can be
improved without improving ranking. Second, error targets are hard to keep. A threshold tuned to 5\% error on hand-coded
comments exceeds 5\% on new comments in about half of re-splits, and applying an independence-based guarantee to clustered
decisions breaks it, because a comment's codes and a transcript's utterances err together. A guarantee that respects how the
hand-coded sample was drawn holds, at a price: after 30\% of Reddit comments are hand-coded, it certifies 60--83\% of the
remaining code decisions for most coders. On dialogue, with whole transcripts hand-coded, it certifies nothing and could not
even for a perfect coder at our calibration sizes; sampling single utterances instead still yields almost no certificates.
Third, accepted shares overstate the work saved: nearly all certified decisions say a code does not apply, and a baseline
that reads no text accepts 75\% of comment decisions at 3.6\% error. Model uncertainty also rises where annotators disagree,
and option order changes some models' labels. Confidence can triage coding, but only if researchers check ranking, sampling
design and what the accepted labels are. Models were called on 10--12 October 2026; we release the harness and every coding
decision.
\end{abstract}
